\begin{table}[htbp]
\centering
\caption{The License-Cap Mechanism: Montana as an Uncapped Contrast Case}
\label{tab:a8_montana}
\small
\begin{tabular}{lcc}
\toprule
Outcome & Coefficient (border effect) & $p$-value (county-cl.) \\
\midrule
\multicolumn{3}{l}{\textit{Colorado + Washington (statewide/local license caps bind)}} \\
log(revenue per store) & $+0.395$ & $0.015$ \\
log(revenue per capita) & $+0.046$ & $0.833$ \\
\midrule
\multicolumn{3}{l}{\textit{Montana (no statewide numerical cap on dispensary licenses)}} \\
log(revenue per capita), 2022--2025 panel & $+0.654$ & $0.036$ \\
log(revenue per store), 2025 cross-section & $+0.012$ & $0.958^{a}$ \\
\bottomrule
\end{tabular}
\begin{minipage}{0.95\textwidth}
\vspace{0.3em}
\footnotesize
\textit{Notes:} $^{a}$HC1 standard error; the Montana revenue-per-store cross-section has too few border
counties ($n=23$) for county-clustered inference. Colorado and Washington cap the number of retail licenses
(Washington explicitly, statewide, at 334 from the outset; Colorado indirectly, through local governments'
authority to ban or cap retail marijuana under Amendment 64). Montana imposes no equivalent statewide cap.
The pattern reverses exactly as the license-cap mechanism predicts: where entry is capped, the border effect
appears in revenue per store and not per capita; where entry is free to expand, it appears in revenue per
capita (more stores absorbing the demand) and not per store.
\end{minipage}
\end{table}
